\majorbreak
\section{Local changes worth trying first}
\label{sec:local}

These edits keep the mathematical interfaces in place. They are deliberately
more specific than advice to replace proofs with \code{simp}, \code{aesop},
or \code{omega}. Each proposed simplification identifies an existing fact
or a repeated invariant. All after blocks remain \textbf{uncompiled}.

\subsection{Specialize the generic coefficient bound already in the file}
\label{sec:halfbase}

\textbf{Priority 1 and the best first pilot.}
\code{SelectiveCoefficientBounds.half_base_bound_31} at
\src{SelectiveCoefficientBounds.lean}{203}{232} reproves a result already
available as \code{half_base_bound} at
\src{SelectiveCoefficientBounds.lean}{166}{199}. The generic theorem has
\code{m} bounded input digits and assumes
\[
 4(m+1)^2Hb^2\le 2^t.
\]
With $m=31$, the scale is exactly $4096Hb^2$, and the specialized theorem's
equality hypothesis supplies the required inequality.

The existing 24-line proof body repeats the sum bound, square bound,
coefficient estimate, strict gap, and predecessor-power calculation. For
example, it ends with this source excerpt:

\begin{Lean}
  have hb2 : 0 < (b : ℤ) ^ 2 := by positivity
  have hgap : 2 * (H * (32 * (b : ℤ)) ^ 2) <
      4096 * H * (b : ℤ) ^ 2 := by
    nlinarith [mul_pos hH hb2]
  have hpow : (2 : ℤ) ^ t = 2 * 2 ^ (t - 1) := by
    rw [← pow_succ']
    congr 1
    omega
  rw [← hbase, hpow] at hgap
  linarith
\end{Lean}

Replace the \emph{whole proof body}, keeping the theorem statement intact:

\begin{Lean}
  refine half_base_bound w u p hH hp hb ht hu ?_ n
  norm_num [hbase]
\end{Lean}

\textbf{Measured candidate comparison:} 24 physical proof-body lines become
2. This is a textual comparison, not an executed proof reduction. No new
API is needed, and the strict conclusion \code{< 2 ^ (t - 1)}, the exact
encoded polynomial, and all positivity/boundedness hypotheses remain.
This is preferable to separately optimizing the duplicated proof's tactics.

\subsection{State the sum of squares invariant directly}
\label{sec:squares}

\textbf{Priority 1.} \code{ShortSystem.residualSum_zero_iff} at
\src{ShortSystem.lean}{63}{88} proves that five integer squares sum to zero
exactly when all five residuals vanish. It currently makes five similar
\code{nlinarith} calls, one for each square. The first is:

\begin{Lean}
  have h0 : (n0 - 1) ^ 2 = 0 := by
    nlinarith [sq_nonneg (n1 - 1), sq_nonneg (na - 1),
      sq_nonneg (ns - 1), sq_nonneg (nk * nt - 1),
      sq_nonneg (n0 - 1)]
\end{Lean}

\existing{} The order API already provides
\code{add_eq_zero_iff_of_nonneg}; the pinned Mathlib source generates this
additive lemma from
\href{https://github.com/leanprover-community/mathlib4/blob/81a5d257c8e410db227a6665ed08f64fea08e997/Mathlib/Algebra/Order/Monoid/Unbundled/Basic.lean\#L1070-L1081}{the multiplicative analogue with both factors at least one}.
Together with \code{sq_eq_zero_iff} and \code{sub_eq_zero}, it describes
the proof's actual invariant.

The concise candidate replaces the complete 23-line body by two lines:

\begin{Lean}
  simp only [residualSum, add_eq_zero_iff_of_nonneg, add_nonneg,
    sq_nonneg, sq_eq_zero_iff, sub_eq_zero, and_assoc]
\end{Lean}

Its main replay uncertainty is whether conditional simplification discharges
the nonnegativity premises as intended in this exact context. A more explicit
fallback proves the generic five-square fact with the partial sums named:

\begin{Lean}
  have hab := add_nonneg (sq_nonneg a) (sq_nonneg b)
  have habc := add_nonneg hab (sq_nonneg c)
  have habcd := add_nonneg habc (sq_nonneg d)
  rw [add_eq_zero_iff_of_nonneg habcd (sq_nonneg e),
    add_eq_zero_iff_of_nonneg habc (sq_nonneg d),
    add_eq_zero_iff_of_nonneg hab (sq_nonneg c),
    add_eq_zero_iff_of_nonneg (sq_nonneg a) (sq_nonneg b)]
  simp only [sq_eq_zero_iff, and_assoc]
\end{Lean}

The archive also gives a direct, explicit source-statement fallback. Choose
one proof, not both. \code{compactResidualSum_zero_iff} already reuses the
original result, so it is not a second independent reduction. Do not import
the heavier \code{Mathlib.NumberTheory.Dioph} module merely for an analogous
sum-of-squares theorem. The lightweight order facts suffice.

\subsection{Use finite set intersection instead of reproving membership}

\textbf{Priority 1 for a small isolated edit.}
\code{LowMaskCoding.sum_range_support} at
\src{LowMaskCoding.lean}{39}{46} first proves that filtering a range by
membership in a contained set gives the set itself. The existing body is:

\begin{Lean}
  have hf : (range len).filter (fun j => j ∈ S) = S := by
    ext j
    simp only [mem_filter]
    exact and_iff_right_of_imp (fun hj => hS hj)
  rw [← Finset.sum_filter, hf]
\end{Lean}

The proposed two-line body uses the existing finite-set API:

\begin{Lean}
  rw [← Finset.sum_filter, Finset.filter_mem_eq_inter,
    Finset.inter_eq_right.mpr hS]
\end{Lean}

\code{Finset.filter_mem_eq_inter} is explicitly present in the
\href{https://github.com/leanprover-community/mathlib4/blob/81a5d257c8e410db227a6665ed08f64fea08e997/Mathlib/Data/Finset/Basic.lean\#L357-L358}{pinned Mathlib source}.
The same change applies to \code{MultiTargetCarry.selectedMask_eq} at
\src{MultiTargetCarry.lean}{55}{67}, converting the pointwise bound
\code{hp i hi} to range membership:

\begin{Lean}
  rw [← Finset.sum_filter, Finset.filter_mem_eq_inter,
    Finset.inter_eq_right.mpr
      (fun i hi => Finset.mem_range.mpr (hp i hi)),
    Finset.mul_sum]
\end{Lean}

The candidate files measure the changed spans as $5\to2$ and $6\to3$
physical lines; the latter is wrapped more generously here for readability.
The mathematical gain is eliminating a second implementation of a standard
set identity, not teaching the simplifier an extra global rule.

\subsection{Remove the repeated predecessor exponent calculation}

\textbf{Priority 2.} Seven local \code{hpow} proofs show
\code{a ^ t = a * a ^ (t - 1)} under \code{1 ≤ t}. They share this
three-line body:

\begin{Lean}
  rw [← pow_succ']
  congr 1
  omega
\end{Lean}

The existing power and natural-subtraction APIs give the proposed body:

\begin{Lean}
  rw [← pow_succ', Nat.sub_add_cancel ht]
\end{Lean}

The occurrences are in
\src{MultiTargetCarry.lean}{40}{43},
\src{MultiTargetCarry.lean}{132}{135},
\src{MultiTargetCarry.lean}{226}{229},
\src{SelectiveCarry.lean}{89}{92},
\src{SelectiveCarry.lean}{111}{114},
\src{SelectiveCoefficientBounds.lean}{194}{197}, and
\src{SelectiveCoefficientBounds.lean}{227}{230}.
One disappears entirely if the 31-variable specialization is adopted, so
its saving must not be counted twice. An extra project lemma is optional;
it is not needed to write this one-line replacement. Keep the \code{ht}
premise: natural subtraction is not cancellative without the bound.

\subsection{Centralize the modular residue interface}
\label{sec:residue}

\textbf{Priority 2 for maintenance rather than raw compression.}
\code{CurrentClause2011.addition_residue} and
\code{subtraction_residue} at
\src{CurrentClause2011.lean}{27}{39} are duplicated privately as
\code{add_residue} and \code{sub_residue} at
\src{CurrentReflectedGeometry2011.lean}{74}{86}.

The arithmetic only needs a divisor \code{L} of the ambient modulus
\code{n}, not the concrete value 2011. A proposed low-level shared module
can expose these exact interfaces:

\begin{Lean}
-- Proposed generic statements; full proofs are in PortableCandidates.lean.
theorem addition_residue {n L : ℕ} [NeZero n] (hL : L ∣ n)
    (b : ZMod n) (e : ℕ) :
    (b + (e : ZMod n)).val % L = (b.val + e) % L := by
  -- Proof omitted here; UNCOMPILED full candidate in the archive.

theorem subtraction_residue {n L : ℕ} [NeZero n] (hL : L ∣ n)
    (k : ZMod n) (e : ℕ) :
    ((k - (e : ZMod n)).val + e) % L = k.val % L := by
  -- Proof omitted here; UNCOMPILED full candidate in the archive.
\end{Lean}

The full candidates use \code{ZMod.natCast_eq_natCast_iff} followed by
\code{Nat.ModEq.of_dvd}, keeping the coercion normalization in one place.
Keep the existing public names as concrete specializations. Do not make
geometry import the clause layer merely to borrow these lemmas, and retain
the positive-modulus instance. The helper-plus-wrapper cost means no large
net line saving is asserted.

\subsection{Smaller opportunities and a reason to keep an existing proof}

The archive contains several further local experiments; they are useful
after the higher-value reuse changes, not a mandate to minimize every proof.

\begin{itemize}
\item \textbf{Convolution factorization.}
In \code{NativeCyclicConvolution.word_convolution}, the final distributivity
step at \src{NativeCyclicConvolution.lean}{79}{87} can be expressed using
\code{Finset.sum_mul}, \code{Finset.mul_sum}, \code{pow_add}, and controlled
associative/commutative simplification. The candidate changes that span from
9 to 3 lines. Preserve the preceding modular-congruence argument; a global
\code{ring} call does not replace it.
\item \textbf{Use the exact power fact.}
At \src{NativeRadixDigits.lean}{54}{57}, replace a weaker power inequality,
normalization, and \code{omega} with
\code{Nat.one_lt_two_pow (by omega)}. The same existing lemma is already
used at \src{NativeWindowCarry.lean}{63}{63}. The candidate changes the
four-line \code{hr} declaration to one line.
\item \textbf{Normalize a finite-switch boundary once.}
At \src{ArithmeticFiniteSwitch.lean}{24}{30}, keep the equality case and
then prove \code{tag < count + 1 ↔ tag < count} in its complement. This
removes the redundant second case split; the supplied candidate span is
7 to 4 lines.
\item \textbf{Do not prefer branch expansion just because it is shorter.}
The optional \code{cellDigit_lt} candidate at
\src{CurrentHighDigits2011.lean}{80}{86} is two lines using
\code{split_ifs} and \code{norm_num}, but duplicates finite/numeral work
already packaged by \code{highDigits_lt}. Keep the original unless actual
readability and profiling favor the replacement.
\end{itemize}

Across these changes, retain explicit \code{simp only} boundaries for large
kernel, stencil, and mask definitions. A project-wide collection of new
\code{[simp]} tags would trade short source for unpredictable expansion in
downstream proofs.
